\section{Verification purpose}

\textbf{Case ID:} \CaseID

This case is a negative verification test for mesh dependence in the current
multiple-discrete-cell treatment of structures. A physical
\SI{10}{m}-by-\SI{10}{m} house may occupy one raster cell, several independent
cells, one oversized cell, or disappear under categorical resampling as the
mesh changes. Because those cells are not assembled into one conserved
structure object, changing resolution alters represented heat release,
firebrand generation, deposition, and structure-to-structure ignition.

The test asks whether all 15 current-fingerprint simulations produce an
evaluable community rate of spread (ROS), and whether the fitted values remain
within the predeclared mesh-independence limits relative to the run-derived
10-m anchor. Detecting the known resolution sensitivity is useful diagnostic
evidence, but it is a \emph{verification failure}, not a pass.

\section{Mathematical formulation and numerical implementation}

\subsection{Mode-resampled square array}

The physical layout is a periodic array of square structures of side
\(a=\SI{10}{m}\), separated in both coordinates by
\(d=\SI{10}{m}\). Define the one-dimensional structural indicator
\[
 I(s)=\begin{cases}1,&s\bmod(a+d)<a,\\0,&\text{otherwise.}\end{cases}
\]
For grid interval \(C_i=[i\Delta,(i+1)\Delta]\), preprocessing calculates
\[
 f_i=\frac{1}{\Delta}\int_{C_i}I(s)\,ds,
 \qquad b_i=\mathbf{1}(f_i\ge 1/2),
 \qquad B_{ij}=b_i b_j .
\]
This coordinate-wise categorical-mode construction reproduces the reported
geometries. At \(\Delta=\SI{7}{m}\), effective side lengths are 7 or
\SI{14}{m} and footprints range from 49 to \(\SI{196}{m^2}\). At
\(\Delta=\SI{10}{m}\), footprint and gap are exact. At
\(\Delta=\SI{30}{m}\), one represented structure cell has area
\SI{900}{m^2} and represented structures are separated by \SI{30}{m}.
The 20- and 40-m maps eliminate represented firebreaks; this is an intended
resampling diagnostic, not a preprocessing failure.

\subsection{Cell-local firebrand treatment}

A represented structural cell uses a design-fire curve \(q^{\prime\prime}(t)\). The curve has a 300-s linear growth phase, a one-hour plateau at \SI{400}{kW.m^{-2}}, and a 300-s decay. Under the current
multiple-discrete-cell treatment,
\[
 \mathrm{HRR}_{ij}(t)=q^{\prime\prime}(t)\Delta^2,
 \qquad \dot N_{ij}(t)=GR\,\frac{\mathrm{HRR}_{ij}(t)}{1000},
 \qquad GR=\SI{10}{pcs.MW^{-1}.s^{-1}}.
\]
Changing \(\Delta\) therefore changes both represented geometry and the HRR
of each independent emitter. Firebrands use the structural empirical
flight-distance distribution, Eulerian accumulation, and physical ignition.
Surface fire is disabled, so every downwind ignition must be caused by
deposited firebrands. The configured time-step scaling is
documented by Qin et al.~\cite{QinEtAl2026}:
\[
 \Delta t=\mathrm{CFL}\frac{\Delta}{U_{\rm wind}},
 \qquad \mathrm{CFL}=1,
 \qquad U_{\rm wind}=\SI{17.9}{m.s^{-1}}\;(\SI{40}{mph}).
\]

For each represented downwind structure column \(s\), postprocessing takes the
median positive finite arrival time \(t_s\). At least three ignited columns are
required. Mean community ROS is the least-squares slope
\[
 ROS_\Delta =
 \frac{\sum_s(t_s-\bar t)(x_s-\bar x)}
      {\sum_s(t_s-\bar t)^2}.
\]
Because the 10-m grid exactly resolves both physical length scales, its fitted
ELMFIRE result is the comparison anchor,
\[
 ROS_{10}=ROS_{\Delta=\SI{10}{m}}.
\]
The dashed reference line and the 10-m ELMFIRE point therefore use the same
calculated value; no rounded external estimate is substituted.

\section{Derivation of expected behavior}
A resolution-independent structure treatment would conserve each house's
physical footprint, integrated heat release, firebrand source, accumulated
load, and ignition state. Subject to numerical error, it would give
\(ROS_\Delta\approx ROS_{10}\) for every resolution and enough downstream
ignitions to fit all 15 ROS values.

The current cell-local representation does not conserve those quantities. Fine
meshes fragment houses into multiple emitters, intermediate meshes alter
footprints and gaps, and some coarse meshes merge structures or remove
firebreaks. The expected diagnostic result is therefore a nonuniform ROS curve,
including possible failure to propagate at some resolutions. That observation
exposes the defect; against the mesh-independence requirement, its correct
overall status is FAIL.

\section{Verification metrics and acceptance criteria}

Preprocessing passes only when all 15 resolutions are generated, the exact
10-m map has no footprint or gap error, the 7-m footprint range is
49--\SI{196}{m^2}, and the median 30-m footprint is \SI{900}{m^2}.
These checks establish that runtime comparison uses the intended
resolution-dependent rasterizations.

A complete runtime decision additionally requires every current-fingerprint
variant and a finite 10-m ROS anchor. Let \(N_{\rm fit}\) be the number of
resolutions with at least three ignited downwind structure columns. ROS
extraction completeness is
\[
 f_{\rm fit}=\frac{N_{\rm fit}}{15},
\]
and must equal one. For the finite fitted values, mesh dependence is measured
by the coefficient of variation
\[
 CV_R=\frac{\operatorname{std}_{\Delta}(ROS_\Delta)}
            {\operatorname{mean}_{\Delta}(ROS_\Delta)}
\]
and by the largest deviation from the exact-geometry anchor,
\[
 D_{\max}=\max_\Delta\frac{|ROS_\Delta-ROS_{10}|}{ROS_{10}}.
\]
The permissive engineering limits are \(CV_R\le0.25\) and
\(D_{\max}\le0.50\). Missing fits are handled by \(f_{\rm fit}\); they are not
silently discarded and counted as convergence. The overall decision is the
conjunction of every geometry, runtime, anchor, fit-completeness, and
mesh-independence check.


\subsection{Justification of metric quantification}

The geometry limits are exact consequences of rasterization.  At
\(\Delta=\SI{10}{m}\), the structure side and gap are both exactly one cell,
so the 1\% footprint and separation allowances cover only file precision and
connected-component bookkeeping.  At \SI{7}{m}, a represented component can
contain one to four cells, giving areas from \(7^2=49\) to
\(4\times7^2=196\ \mathrm{m^2}\).  At \SI{30}{m}, a one-cell represented
structure has area \(30^2=900\ \mathrm{m^2}\).  These checks establish that
the intended mesh-dependent maps were generated; they are not claims of
geometric accuracy.

For the behavioral requirement, \(f_{\rm fit}=1\) is non-negotiable because
mesh independence cannot be inferred after silently discarding resolutions
that fail to propagate.  The \(CV_R\le0.25\) and \(D_{\max}\le0.50\)
limits are intentionally permissive engineering bounds, substantially larger
than ordinary front-fit quantization.  Even with those allowances, a method
with factor-of-two variation or missing fits fails, which is the defect this
case is designed to expose~\cite{QinEtAl2026}.  The 10-m reference is the
actual \(ROS_{10}\) extracted by the same postprocessor, so its plotted
reference and ELMFIRE point must coincide by construction; it is not an
independent nominal value.  Current fingerprints and all 15 runs are required
to prevent stale or selectively complete outputs from weakening the expected
FAIL diagnosis.

\subsection{Predeclared acceptance contract}

Table~\ref{tab:acceptance-contract} summarizes the metrics fixed before
inspection of the calculated results. Numerical values and component statuses
are reported separately in Section~7.

\begin{table}[H]
\centering
\footnotesize
\caption{Predeclared verification metrics and acceptance criteria.}
\label{tab:acceptance-contract}
\begin{tabular}{@{}p{0.23\linewidth}p{0.47\linewidth}p{0.21\linewidth}@{}}
\toprule
Metric & Output, region, and aggregation & Acceptance criterion\\
\midrule
Prepared variants & Generated manifest count divided by the 15 requested grid spacings. & 100\% \\
10-m footprint error & Maximum relative error of represented component areas from \SI{100}{m^2}. & \(\le0.01\) \\
10-m separation error & Maximum relative gap-width error from \SI{10}{m}. & \(\le0.01\) \\
7-m component areas & Minimum and maximum represented connected-component areas. & 49--\SI{196}{m^2} \\
30-m footprint error & Relative error of median represented footprint from \SI{900}{m^2}. & \(\le0.01\) \\
Runtime completeness & Current-fingerprint final TOA output and successful termination. & 15 of 15 variants \\
10-m ROS anchor & Structure-column fit from the current 10-m TOA raster. & Finite \(ROS_{10}\) \\
ROS extraction completeness & Resolutions with at least three ignited downwind columns. & \(f_{\rm fit}=1\) \\
ROS mesh-dependence CV & Coefficient of variation over finite fitted ROS values. & \(CV_R\le0.25\) \\
Maximum anchor deviation & Maximum finite-grid relative deviation from \(ROS_{10}\). & \(D_{\max}\le0.50\) \\
Overall decision & Every geometry, completeness, and ROS component is conjunctive. & PASS only if all pass \\
\bottomrule
\end{tabular}
\end{table}

\section{Simulation configuration and predicted outcome}

Figure~\ref{fig:input-configuration} records the prepared spatial inputs for \texttt{dx10}.
The postprocessor reads these panels from the variant-local GeoTIFF files, removes the
two-cell numerical buffer, and retains map coordinates in metres.  The left panel shows
the most informative fuel or structure layout available for the case; the right panel
shows the cells selected by the non-positive initial level-set field, making a localized
ignition or firebrand-emission source explicit.

\begin{figure}[H]
\centering
\includegraphics[width=0.98\textwidth,height=0.42\textheight,keepaspectratio]{../figures/input_configuration.pdf}
\caption{Whole-domain prepared input configuration for the representative variant.  The
physical domain is shown without the numerical halo; coordinates, configured wind values,
fuel or structure layout, and initial ignition/source geometry are identified in the panels.}
\label{fig:input-configuration}
\end{figure}


\begin{table}[H]
\centering\footnotesize
\caption{Configuration of the two-dimensional resolution sweep.}
\begin{tabular}{@{}p{0.22\linewidth}p{0.27\linewidth}p{0.13\linewidth}p{0.29\linewidth}@{}}
\toprule
Parameter & Configured value & Units & Role\\
\midrule
Physical domain & \(1200\times400\), flat & m & permits many downwind ignitions\\
Structure side / gap & 10 / 10 & m & exact physical scales at the 10-m grid\\
Initial ignition & complete first structure column & -- & uniform upwind source across all structure rows\\
Grid spacings & 1, 3, 5, 7, 10, 13, 15, 17, 20, 23, 25, 27, 30, 40, 50 & m & exposes raster and cell-local dependence\\
Numerical halo & 2 on every side & cells & excluded from physical geometry\\
Wind & 40 from 270 degrees & mph & constant west-to-east transport\\
Time step & \(\Delta/U_{\rm wind}\) & s & holds wind CFL at one\\
Design-fire peak & 400 & \si{kW.m^{-2}} & structural HRRPUA input\\
Generation rate & 10 & \si{pcs.MW^{-1}.s^{-1}} & cell-local building source\\
Transport / accumulation & empirical / Eulerian & -- & current structural spotting path\\
Ignition delays & 42 plus 300 & s & small-flame and development stages\\
Surface spread & disabled & -- & isolates firebrand-driven propagation\\
Stop time & 150000 & s & permits slow community spread\\
Final output & one TOA per variant & -- & source for structure-column ROS fits\\
\bottomrule
\end{tabular}
\end{table}

The dashed reference line is predicted to intersect the 10-m ELMFIRE point
exactly because both are \(ROS_{10}\). A mesh-independent implementation would
cluster all other finite ROS points around this line and provide a fit at every
resolution. The current independent-cell treatment is instead expected to
produce missing fits and large variations, causing the behavioral verification
to fail.

\subsection{Preprocessing, execution, and postprocessing}

\texttt{scripts/preprocess.py} places customizable constants at its beginning,
constructs the periodic map analytically, assigns persistent structure IDs,
writes co-registered GeoTIFF inputs and case-local fuel tables, and records a
SHA-256 fingerprint for every variant. It also writes
\texttt{outputs/geometry\_summary.csv}, making the map checks independently
auditable.

\texttt{run\_case.sh} is executable within the case directory and follows the same linear workflow as the other manifest cases. Configure \texttt{ELMFIRE\_BIN} before launching it; the runner executes all 15 variants sequentially, keeps logs and fingerprints variant-local, postprocesses the outputs, and compiles the report. Because the sweep is long, execution should be scheduled only when adequate compute time is available.

\texttt{scripts/postprocess.py} rejects stale or incomplete outputs, reads the
final time-of-arrival raster, groups cells by structure, fits ROS, calculates
all metrics above, and produces the report figure.

\section{Actual simulation results}

Figure~\ref{fig:domain-result} shows the representative time-of-arrival field from \texttt{dx10}.
It is read from the actual ELMFIRE output named in the panel, masks nodata, and excludes
the two-cell numerical buffer.  The 10-m field is the reference-scale spatial result used to interpret mesh-dependent structure spread across the resolution sweep.

\begin{figure}[H]
\centering
\includegraphics[width=0.96\textwidth,height=0.50\textheight,keepaspectratio]{../figures/domain_result.pdf}
\caption{Whole-domain time-of-arrival field from the representative completed ELMFIRE run.  This
domain-scale result supplements, rather than replaces, the higher-level verification
profiles, convergence plots, ensemble statistics, and calculated acceptance metrics below.}
\label{fig:domain-result}
\end{figure}


All \CompletedVariantCount{} of \RequestedVariantCount{} current-fingerprint
variants completed successfully. The selected output for each variant is
\path{time_of_arrival_0000001_0150000.tif}. The overall decision
calculated by postprocessing is \textbf{\OverallStatus}.

The upper row of Figure~\ref{fig:spatial-resolution} shows the actual
preprocessed geometries. Its lower panel contains fitted ELMFIRE ROS values.
Red crosses at zero identify completed simulations with fewer than three
ignited downwind columns, for which ROS is deliberately not fitted. The dashed
line is \(ROS_{10}=\DxTenReferenceROS~\si{m.s^{-1}}\), exactly equal to the
10-m ELMFIRE point.

\begin{figure}[H]
\centering
\includegraphics[width=\linewidth]{spatial_resolution_sweep.pdf}
\caption{Two-dimensional spatial-resolution diagnostic. Top: mode-resampled
WUI geometry at 7, 10, and \SI{30}{m}. Bottom: mean community ROS from current
ELMFIRE TOA outputs; red crosses identify resolutions without enough downwind
ignitions for a fit, and the dashed line is the run-derived 10-m anchor. The
geometry selection follows the documented comparison
configuration~\cite{QinEtAl2026}.}
\label{fig:spatial-resolution}
\end{figure}

\section{Calculated metrics and verification decision}

\begin{table}[H]
\centering\small
\caption{Calculated verification metrics and component decisions.}
\begin{tabular}{@{}p{0.37\linewidth}p{0.18\linewidth}p{0.17\linewidth}p{0.14\linewidth}@{}}
\toprule
Metric & Acceptance limit & Calculated & Status\\
\midrule
\MetricRows
\bottomrule
\end{tabular}
\end{table}

The postprocessor reports \textbf{\OverallStatus}. Geometry preparation and
all 15 executions pass their completeness checks, and the run-derived 10-m
anchor is finite. The behavioral checks fail: the ROS-evaluable fraction is
\ROSEvaluableFraction, the finite-grid coefficient of variation is
\ROSMeshCV, and the maximum relative deviation from the 10-m anchor is
\ROSMaximumDeviation.

\section{Technical interpretation}
Only seven of the 15 completed simulations ignite enough downwind structure
columns to support a ROS fit. Among the finite fits,
\(ROS{10}=\DxTenReferenceROS~\si{m.s^{-1}}\), while the variation is large:
\(CV_R=\ROSMeshCV\) and \(D_{\max}=\ROSMaximumDeviation\). The 40-m run,
for example, reaches approximately \SI{0.04054}{m.s^{-1}}, more than three
times the 10-m anchor.

This combination of missing propagation at eight resolutions and widely
varying finite ROS values is direct evidence that the current independent-cell
house representation is mesh dependent. The FAIL status is therefore the
intended diagnostic interpretation. It must not be converted to PASS merely
because the observed nonconvergence agrees with a known limitation.

\section{Scope and limitations}
This case verifies numerical mesh dependence for one idealized periodic WUI
array, one wind, and one structural design fire. It does not validate the
physical accuracy of any individual ROS value. The 10-m result is an internal
run-derived anchor, not experimental validation data. A future
resolution-independent structure implementation should be tested with the same
geometry and acceptance contract; it should provide all 15 ROS fits and reduce
both variation metrics below their limits.

\begin{thebibliography}{9}
\bibitem{Qin2025}
Y. Qin, \emph{A Physics-Based Eulerian Framework for Modeling Firebrand
Showering in Regional-Scale Wildland and Wildland--Urban Interface Fire
Simulations}, Ph.D. dissertation, University of Maryland, College Park, 2025.

\bibitem{QinEtAl2026}
Y. Qin et al., ``Simulations of firebrand-driven fire spread in
landscape-scale Wildland--Urban-Interface and urban conflagration models,''
\emph{Fire Safety Journal}, vol.~162, 104686, 2026,
doi:10.1016/j.firesaf.2026.104686.
\end{thebibliography}
